\documentclass[11pt]{article}
\usepackage[utf8]{inputenc}
\usepackage{amsmath,amssymb,amsthm}
\usepackage[margin=1in]{geometry}
\usepackage{hyperref}
\usepackage{xcolor}

\newcommand{\tideref}[2][]{\href{https://therisensea.org/tides/#2/}{\texttt{#2}}}

\title{Tide retrospective: the expansion pairing reads off the
quartic coefficient}
\author{Tide loop (Claude Fable 5, with GPT-5.6 Sol deliberation)}
\date{2026-08-09}

\begin{document}
\maketitle

\section{Setting}

The germbij note's Section~7.4 recovers the subleading coefficients of an
analytic potential from the corrections of the Laplace expansion at their
distinct orders. The recovery thread had so far certified the leading
data (degree and scale, exactly and asymptotically) and a monotonicity
shortcut for signed perturbations; the expansion \emph{mechanism} itself,
the pairing the note actually describes, had no Lean instance. This tide
provides the smallest one: the first correction of the quartic-perturbed
Gaussian, with an explicit error bound, and the recovery of the quartic
coefficient from it.

\section{The thing formalised}

Three declarations in \texttt{Laplace/OneD/RecoveryExpansion.lean},
sorry-free with zero warnings. The elementary remainder
bound\footnote{\texttt{exp\_neg\_sub\_one\_add\_bounds}.}: for $s \geq 0$,
$0 \leq e^{-s} - 1 + s \leq s^2$, proved with two applications of
$1 + x \leq e^x$ and no calculus. The
expansion\footnote{\texttt{quartic\_partition\_expansion\_bounds}.}: for
$b \geq 0$, $t > 0$,
\[
0 \;\leq\; Z_b(t) - \sqrt{2\pi}\, t^{-1/2} + 3b\sqrt{2\pi}\, t^{-3/2}
\;\leq\; 105\, b^2 \sqrt{2\pi}\, t^{-5/2},
\qquad Z_b(t) = \int e^{-t(x^2/2 + b x^4)}\, dx .
\]
And the recovery\footnote{\texttt{quartic\_coefficient\_recovery\_of\_eventuallyEq}.}:
if $Z_{b_1}$ and $Z_{b_2}$ eventually agree ($b_i \geq 0$), then
$b_1 = b_2$. The numerical check (executed, archived in the tide log)
confirms both the limit $-3b\sqrt{2\pi}$ of the rescaled correction and
the error constant with the expected factor-two slack from the
elementary bound.

\section{Proof strategy}

Split $e^{-t(x^2/2+bx^4)} = q \cdot e^{-tbx^4}$ with $q$ the Gaussian
factor, and expand $e^{-tbx^4} = 1 - tbx^4 + r$ with
$0 \leq r \leq (tbx^4)^2$ pointwise. The seabed's scale-$t$ moments (the
$k = 0, 2, 4$ cases of $\int x^{2k} q = (2k-1)!!\,\sqrt{2\pi}\,
t^{-(k+1/2)}$, so $1$, $3$, $105$) do the rest: the $k = 0$ and $k = 2$
terms cancel the main and correction terms of the expansion
\emph{exactly}, so the error is literally $\int q\, r$, bounded two-sidedly
by the $k = 4$ moment. No substitution, no dominated-convergence, no
asymptotic machinery. For the recovery, the difference of the two
expansion errors is exactly $3(b_1 - b_2)\sqrt{2\pi}\, t^{-3/2}$ when the
partition functions agree, and the two error bounds squeeze it below
order $t^{-5/2}$; one archimedean choice of $t$ closes the contradiction,
with no limits taken.

\section{Roadblocks and resolutions}

Four build iterations, all elaboration-level, all of documented classes:
a missing import for the seabed integrability family; the
\texttt{Pi.add} gotcha in \texttt{integral\_add} (fixed by the
CLAUDE.md's type-ascribed single-lambda witnesses); brittle
\texttt{rw}-into-closed-form steps replaced by the robust pattern
``\texttt{norm\_num [Nat.doubleFactorial]} at hypothesis and goal, then
\texttt{linarith} with the integral as an atom''; and a \texttt{calc} of
equalities against a $\leq$ goal needing \texttt{le\_of\_eq}.

\section{What was Mathlib, what was new}

Mathlib: $1 + x \leq e^x$, integral linearity, \texttt{integral\_mono},
domination. Seabed: the scale-$t$ Gaussian moments and the monomial
integrability family, both from the May tides, doing exactly the work
they were built for. New: the observation that the moment cancellation
is exact (the expansion error \emph{is} the remainder integral), which
makes the whole development bookkeeping-free.

\section{Lessons}

When a closed form is consumed, normalize both sides to a common normal
form and finish with \texttt{linarith} treating opaque terms as atoms;
chasing exact \texttt{rw} orientations against \texttt{norm\_num}
output is the fragile alternative. The elementary $s^2$ remainder bound
(constant twice the sharp $s^2/2$) is the right trade: it removes all
calculus from the file at the cost of a factor the application never
notices.

\section{Follow-ups}

\begin{itemize}
\item The general inductive step (all subleading coefficients of an
analytic perturbation, correction by correction) would need expansions
to arbitrary order: the natural shape is a finite Taylor expansion of
$e^{-s}$ with the same elementary remainder trick, and it would subsume
this tide's $k = 1$ instance.
\item Tag the note's Section 7.4 dimension-one item with the expansion
recovery once merged; it certifies the pairing mechanism specifically.
\end{itemize}

\end{document}
